\documentclass[11pt]{article}
\usepackage[utf8]{inputenc}
\usepackage[T1]{fontenc}
\usepackage{lmodern}
\usepackage{textcomp}
\usepackage{amsmath, amssymb, amsthm}
\usepackage{graphicx}
\usepackage{listings}
\usepackage[hidelinks]{hyperref}

\theoremstyle{plain}
\newtheorem{theorem}{Teorema}[section]
\newtheorem{lemma}[theorem]{Lema}
\newtheorem{proposition}[theorem]{Proposición}
\newtheorem{corollary}[theorem]{Corolario}

\theoremstyle{definition}
\newtheorem{definition}[theorem]{Definición}


\title{c8}
\date{}
\begin{document}
\maketitle

\section{Chapter 8}

\section{Induction}

Induction is a fundamental method of proof, used in every corner of mathematics. It is applicable to statements of the form

$\forall$n $\in$N, P(n) (8.1)

where P is a predicate over N. This sentence says that all natural numbers have property P.

Induction is first met in the proof of summation identities and inequalities:

n

xk = 1 $-$xn+1

x = 1 Sum of geometric progression

1

1 $-$x

k=0

n

n

2 (x + y)n =

xk yn$-$k Binomial theorem

k

k=0

3 (1 + x)n $\geq$1 + nx x > $-$1 Bernoulli inequality1

n

n

n

1 n

xk > 0 Arithmetico-geometric

xk $\leq$

xk

4

k=1

k=1

mean (AGM) inequality

2

n

n

n

|xk|2

|yk|2 Cauchy-Schwarz inequality2

xk yk

5

$\leq$

k=1

k=1

k=1

(The variables are complex numbers in 1, 2, and real numbers in 3, 4, 5.)

The expression (8.1) unfolds into an infinite sequence of boolean expressions:

P(1), P(2), P(3), . . .

1 Jacob Bernoulli (Swiss: 1654--1705). 2 Augustin-Louis Cauchy (French: 1789--1857); Hermann Schwarz (German: 1843--1921).

© Springer-Verlag London 2014 F. Vivaldi, Mathematical Writing, Springer Undergraduate Mathematics Series, DOI 10.1007/978-1-4471-6527-9\_8

139

140 8 Induction

and to prove it we must show that all expressions in the sequence are true:

T, T, T, \dots{}

For instance, the Bernoulli inequality unfolds as follows:

1 + x $\geq$1 + x, (1 + x)2 $\geq$1 + 2x, (1 + x)3 $\geq$1 + 3x, . . .

each inequality being valid for all real numbers x > $-$1.

A special form of mathematical argument is needed for this purpose, called mathematical induction. It takes several forms, and we will examine four:

(A) The well-ordering principle. (B) The infinite descent method. (C) The induction principle.

(D) The strong induction principle.

We will show that these principles are equivalent in the sense that any one of them implies all the others. The different forms of the principle are appropriate for different situations.

The term mathematical induction was coined by De Morgan around 1840. The attribute `mathematical' is used to differentiate this concept from the homonymous concept in philosophy, which has a different meaning.3 The induction principle---in any of the above forms---is one of the five Peano axioms4 which characterise the set N of natural numbers.

\section{8.1 The Well-Ordering Principle}

\section{The well-ordering (or least counter example) principle states that}

(A) Every non-empty set of natural numbers contains a least element.

This principle is clearly false for subsets of Z, Q or R, and, less trivially so, for subsets of Q+ or R+, or of any real or rational interval. A possible strategy for proving (8.1) is to combine well-ordering with contradiction. Such a proof will begin as follows:

3 In philosophy, an inductive argument is one that uses very strong premises to support the probable truth of a conclusion. 4 Giuseppe Peano (Italian: 1858--1932).

8.1 The Well-Ordering Principle 141

\begin{proof}
Proof. Suppose (8.1) is false. Let k be the least natural number n for which P(n) is false. \dots{}
\end{proof}

Then with this k we try to deduce a contradiction, using the knowledge that P(n) is true for n = 1, . . . , k $-$1. The integer k defined in the proof exists by virtue of the well-ordering principle: it is the least element of the set \{n $\in$N : $\neg$P(n)\}, which is a subset of N. By assumption, this subset is non-empty.

Let us apply this strategy to some specific problems.

Proposition. If n $\geq$7, then n! > 3n.

\begin{proof}
Proof. Suppose this statement is false, and let k be the least integer n $\geq$7 such that n! $\leq$3n. We know that k > 7, since
\end{proof}

7! = 5040 > 2187 = 37.

Put j = k $-$1. Then j $\geq$7, and hence j! > 3 j by choice of k. So

k! = ( j + 1) j! > ( j + 1)3 j > 3 \textperiodcentered{} 3 j = 3k

contradicting the choice of k. 

The initial verification that k > 7 creates enough room for the principle to work. We note that 6! = 720 and 36 = 729, so the inequality n $\geq$7 is the best possible one. Under these circumstances, we say that n $\geq$7 is a strict bound. Our next statement is more substantial:

Theorem. Every integer greater than 1 is a product of one or more prime numbers.

\begin{proof}
Proof. Suppose the statement is false, and let k be the smallest integer greater than 1 which is not a product of primes. Then k is not prime, so k = mn for two smaller integers m, n $\geq$2. Since m, n are smaller than k and greater than 1, they are products of prime numbers. So k = mn is also a product of primes, which contradicts the choice of k. 
\end{proof}

This form of mathematical induction can be very effective. The assumption that k is the smallest counterexample puts the maximum amount of information on the table; it gives us the feeling that we are examining a concrete object.

\section{8.2 The Infinite Descent Method}

This is a second form of mathematical induction, developed in the 17th Century by Fermat.5 It exploits the following variant of the well-ordering principle:

(B) There is no infinite decreasing sequence of natural numbers.

5 Pierre de Fermat (French: 1601 or 1607/08--1665).

142 8 Induction

This principle follows from the well-ordering principle ((A) $\Rightarrow$(B)): if there were an infinite decreasing sequence of natural numbers

n1 > n2 > n3 > \textperiodcentered{} \textperiodcentered{} \textperiodcentered{}

then the set \{n1, n2, n3, . . .\}, which is a non-empty subset of N, would have no smallest element. Conversely, well-ordering follows from descent ((B) $\Rightarrow$(A)): if there were a non-empty set S $\subset$N without a smallest element, then choosing an arbitrary element n1 $\in$S, we could find n2 $\in$S with n2 < n1. Repeating this process we would generate an infinite decreasing sequence of natural numbers.

Fermat applied this method extensively to arithmetical problems. He pointed out that infinite descent is a method for proving that certain properties or relations for whole numbers are impossible, and for this purpose he combined descent with contradiction---as we did for well-ordering. To prove that no positive integer can have a certain set of properties, we assume that a positive integer does have these properties, and we try to deduce that there is a smaller positive integer with the same set of properties. If we succeed, then, by the same argument, these properties would hold for some smaller positive integer, and so forth ad infinitum. Because a sequence of natural numbers cannot decrease indefinitely, no positive integer can have the stated properties.

Infinite descent was brought to prominence by Euler, who proved with it that it is impossible to find natural numbers x, y, z such that

x3 + y3 = z3.

This is a special case of the celebrated Fermat's last theorem, in which the exponent 3 is replaced by an arbitrary integer n > 2. In keeping with this tradition, we now apply the method of infinite descent to prove the non-existence of integer solutions of a polynomial equation.

Theorem. For every prime p the equation px2 = y2 has no solutions in the natural numbers.

This statement is equivalent to the irrationality of $\surd$p, as one sees by dividing each term of the equation by x2, and then taking the positive square root. The argument presented here is a variant and a generalisation of that given in Sect. 7.1 for the irrationality of

$\surd$

2.

\begin{proof}
Proof. Suppose that there are natural numbers n0, n1, such that pn2
\end{proof}

1 = n2 0. So pn2 1 is the square of a natural number. Since p divides n2

0 = n0n0, by the basic property of prime numbers p must divide n0, and we have n0 = pn2 for some natural number n2. Then pn2

1 = p2n2 2, and hence n2 1 = pn2 2. So n2 < n1 and pn2 2 is the square of a natural number.

Repeating the argument, there is n3 < n2 such that pn2

3 is the square of a natural number. This continues for ever, which is impossible. Hence the equation px2 = y2

has no solution, as stated. 

8.3 Peano's Induction Principle 143

\section{8.3 Peano's Induction Principle}

This is a third form of mathematical induction; it was proposed by Dedekind6 and formalised by Peano. The principle of induction states that:

(C) If P is a predicate over N such that

i) P(1) is true; ii) $\forall$k $\in$N, P(k) $\Rightarrow$P(k + 1);

then P(n) is true for all n $\in$N.

The first condition is called the basis of the induction (or the base case); the second the inductive step. The bi-unique correspondence between predicates P over N and subsets S of N, given by [see Eqs. (4.15 and 4.16)]

S = \{n $\in$N : PS(n)\}, PS(x) = (x $\in$S)

allows us to reformulate the induction principle in the language of sets.

(C$'$) If S is a subset of N such that

i) 1 $\in$S; ii) $\forall$k $\in$N, k $\in$S $\Rightarrow$(k + 1) $\in$S;

then S = N.

The principle of induction follows from well-ordering ((A) $\Rightarrow$(C)). To show this, we define the set

S $'$ = N$\setminus$S = \{n $\in$N : $\neg$P(n)\}.

We prove that S $'$ is empty by contradiction. If S $'$ is non-empty, then, by the well-ordering axiom, S $'$ has a least element k. Since 1 $\in$S, we have k = 1, and since k is the smallest element of S $'$, it follows that k $-$1 $\in$S. Then, putting j = k $-$1, we find from condition ii) that j + 1 = k $\in$S, which contradicts the fact that k $\in$S $'$.

The base case of Peano's induction could be any integer, not just 1. Indeed if P(n) is valid for all n $\geq$k, then by letting m = n $-$k +1, the range n $\geq$k becomes m $\geq$1.

Peano's induction provides straightforward proofs of finite sums and products formulae:

n

a j = F(n) (8.2)

$\forall$n $\in$N,

j=1

n

$\forall$n $\in$N,

a j = G(n) (8.3)

j=1

6 Richard Dedekind (German: 1831--1916).

144 8 Induction

where (a j) is a sequence of numbers (or more generally of elements of a commutative ring), and F and G are explicit functions of n, hopefully easier to compute than the original sum or product.

In an inductive proof of (8.2), the base case consists of verifying that a1 = F(1). For the inductive step, we assume that the formula holds for some n = k, and using the induction hypothesis we obtain:

k+1

k

a j =

a j + ak+1 = F(k) + ak+1.

j=1

j=1

Thus the proof reduces to the verification of the statements

a1 = F(1), F(k) + ak+1 = F(k + 1) k $\geq$1 (8.4)

which no longer involve summation. For example, to prove the formula

n

j2 = n(n + 1)(2n + 1)

n = 1, 2, . . . 6

j=1

we must verify that

11 = 1 \textperiodcentered{} 2 \textperiodcentered{} 3 k(k + 1)(2k + 1)

+ (k + 1)2 = (k + 1)(k + 2)(2k + 3) 6

,

. 6

6

For a product formula, the expressions (8.4) are replaced by

a1 = G(1), G(k) $\times$ ak+1 = G(k + 1) k $\geq$1. (8.5)

In some cases the identities (8.4) and (8.5) may be established using a computer algebra system. This is an instance of computer-assisted proof7---see Exercise 8.1.

A similar arrangement applies to the inductive proof of inequalities. However, even in the simplest cases, such proofs are not as mechanical as those of summation and product formulae, as the following example illustrates.

Proposition. If n > 3, then 2n $\geq$n2.

\begin{proof}
Proof. We prove it by induction on n. The base case n = 4 is immediate: 24 $\geq$42. Assume now that for some k $\geq$4 we have 2k $\geq$k2. Then
\end{proof}

2k+1 = 2 \textperiodcentered{} 2k $\geq$2 \textperiodcentered{} k2

where the inequality follows from the induction hypothesis. To complete the proof, we must show that for all k $\geq$4 we have 2 \textperiodcentered{} k2 $\geq$(k + 1)2. Now the polynomial

7 This concept raises controversy among mathematicians.

8.3 Peano's Induction Principle 145

p(k) = 2 \textperiodcentered{} k2 $-$(k + 1)2 = k2 $-$2k $-$1

$\surd$

has roots 1 $\pm$

2, with

$\surd$ 2 < 1 +

$\surd$

2 < 3.

$-$1 < 1 $-$

Thus, for k $\geq$3, we have p(k) > 0, or 2k2 > (k + 1)2. Hence 2k+1 $\geq$(k + 1)2, completing the induction. 

Thenextexamplewarnsusthatastraightforwardinductivestrategydoesn'talways work.

Proposition. For all integers n $\geq$0 and real numbers x > $-$1, we have

(1 + x)n > nx. (8.6)

We attempt to prove this by induction on n.

Base case: n = 0. Since x > $-$1, we have (1 + x)0 = 1 > 0 = 0x, as desired. Inductive step: Assume (8.6) for some n = k $\geq$0. Then

(1 + x)k+1 = (1 + x)(1 + x)k

> (1 + x)kx (by induction hypothesis) = kx + kx2 ???

We don't seem to be getting anywhere.

\begin{proof}
Proof. We prove the stronger inequality (1 + x)n $\geq$1 + nx. This is the Bernoulli inequality---see p. 139. Base case: n = 0. Since x > $-$1, we have (1 + x)0 = 1 $\geq$1 + 0x. Inductive step: Assume the inequality for n = k $\geq$0. Then
\end{proof}

(1 + x)k+1 = (1 + x)(1 + x)k

$\geq$(1 + x)(1 + kx) (by induction hypothesis) = 1 + (k + 1)x + kx2 $\geq$1 + (k + 1)x.

The inductive step is complete. 

\section{8.4 Strong Induction}

In an induction proof, the successful completion of the inductive step rests on the assumption that to prove P(k + 1) we only need P(k). Should we instead (or in addition) need P(k $-$1), then this method would collapse.

This situation is not unusual. Suppose we want to prove by induction that every integer is a product of prime numbers. To prove that this statement is true for, say,

146 8 Induction

k + 1 = 12, the preceding case k = 11 is of no use. We would need instead k = 3 and k = 4.

The rescue is a form of induction called strong induction (or complete induction), which is formulated as follows:

(D) If P is a predicate over N such that

i) P(1) is true; ii) for any k $\in$N, if P(1), P(2), . . . , P(k) are true, so is P(k + 1);

then P(n) is true for all n $\in$N.

As with Peano's induction, the base case of strong induction can be any integer, and there is an equivalent formulation in terms of sets:

(D$'$) If S is a subset of N such that

i) 1 $\in$S; ii) for any k $\in$N, \{1, 2, . . . , k\} $\subset$S $\Rightarrow$k + 1 $\in$S;

then S = N.

Strong induction follows from Peano's induction ((C) $\Rightarrow$(D)), since the latter has fewer conditions. (One could say that induction is stronger than strong induction!)

In Sect. 8.1 we proved that every integer greater than 1 is a product of one or more primes by combining well-ordering with contradiction; we now prove the same statement using strong induction. The two proofs employ two variants of the same argument, giving a tangible illustration of the equivalence of well-ordering and strong induction.

\begin{proof}
Proof. We use strong induction. The integer 2 is prime, which establishes the base case. Suppose that for some k > 1 every integer n with 1 < n $\leq$k is a product of primes, and consider k + 1. If k + 1 is prime, then there's nothing to prove. If k + 1 is not a prime then there are integers m, n such that mn = k + 1 and 1 < m, n $\leq$k.Bytheinductionhypothesis,m andn areproductsofprimes,andhence so is k + 1. We have completed the inductive step. 
\end{proof}

With the help of strong induction, Euclid's elegant proof of the infinitude of primes (Sect.7.5) yields an upper bound (or an estimate) for the nth prime number.

\begin{theorem}
Theorem 8.1 For all n $\geq$1, the nth prime pn satisfies the upper bound pn $\leq$22n$-$1.
\end{theorem}

\begin{proof}
Proof. We use strong induction. The base case is clear: p1 = 2 $\leq$220 = 2. Assuming that for some k $\geq$1 we have p j $\leq$22 j$-$1 for j = 1, 2, . . . k, we have
\end{proof}

8.4 Strong Induction 147

k

k

22 j$-$1

pk+1 $\leq$1 +

p j $\leq$1 +

j=1

j=1

k

2 j$-$1

= 1 + 22k$-$1 < 22k

j=1

= 1 + 2

where the first inequality follows from Euclid's argument (the right-hand side is divisible by a prime greater than pk, see (7.7)), while the second inequality follows from the induction hypothesis. The proof is complete. 

To complete the proof that the four formulations of the induction principle are equivalent, it remains to show that well-ordering follows from strong induction ((D) $\Rightarrow$(A)). When this is done, we will have proved that

(A) $\Leftrightarrow$(B) and (A) $\Rightarrow$(C) $\Rightarrow$(D) $\Rightarrow$(A)

and equivalence follows from a loop of implications (cf. Sect. 7.4.1).

We use the both ends method. Assume that strong induction holds but that there is a non-empty subset S of N without a least element. Let S $'$ = N$\setminus$S. Then 1 $\in$S $'$,

for otherwise 1 would be the least element of S. Likewise, if for some k $\geq$1 we have \{1, . . . , k\} $\subset$S $'$, then we must also have k + 1 $\in$S $'$, because otherwise k + 1 would be the smallest element of S. Hence, from strong induction, we have that S $'$ = N,

and hence S = $\emptyset$, contradicting the hypothesis S = $\emptyset$.

\section{8.5 Good Manners with Induction Proofs}

Induction proofs are easy to structure, as there is a clear procedure to follow. But even if the steps in a proof are predictable, they should be spelled out clearly.

1. At the beginning of the proof, say with respect to what variable induction is performed, unless this is obvious.

We proceed by induction on the degree d of the polynomial.

2. If the induction variable is n, then in the inductive step it is considered good practice to go from n = k to n = k + 1 rather than from n to n + 1.

3. In the inductive step, say at what point the induction hypothesis is used.

\dots{}where the last inequality follows from the induction hypothesis.

4. Say when the inductive step is complete.

This completes the induction.

If an induction argument is easy, then it may be tempting to skip it altogether.

n, for n $\geq$1. Then xn = 22n$-$1.

Let x1 = 2, and let xn+1 = x2

148 8 Induction

Thelastsentenceisabitblunt,anditcouldbereplacedbymoreconsideratesentences, such as:

A straightforward induction argument shows that xn = 22n$-$1.

Then x2 = 22, x3 = 222, and, in general, xn = 22n$-$1.

Exercise 8.1 Prove the following formulae by induction, first analytically, then with a computer-assisted proof.

n

(2k $-$1) = n2

1.

k=1

n

2.

k! \textperiodcentered{} k = (n + 1)! $-$1

k=1

n

2

n

k3 =

3.

k

k=1

k=1

n

1 n k(k + 1) = n + 1

4.

k=1

n

2n + 1 4 $-$ 2n(n + 1)

k2 $-$1 = 3 1

5.

k=2

n$-$1

(1 + x2k) = 1 $-$x2n

x = 1. 1 $-$x

6.

k=0

Exercise 8.2 Prove by induction all identities and inequalities listed at the beginning of this chapter. (The Bernoulli inequality is proved in Sect. 8.3.)

Exercise 8.3 Prove the following statements by induction.

1. The sum of the cubes of three consecutive natural numbers is a multiple of 9.

2. The sum of the internal angles of a polygon with n sides is (n $-$2)$\pi$.

3. The negation of a boolean expression with n leading quantifiers is obtained by

interchanging quantifiers and negating the predicate.

4. The regions in which the plane is subdivided by n lines may be coloured with just

two colours in such a way that no two neighbouring regions have same colour.

Exercise 8.4 Find each formula, and then prove it by induction.

1. The formula for the cardinality of the power set of a finite set.

2. The formula for the higher-order derivatives of the product of two functions.

3. The formula for the derivative of the repeated composition of a function with

itself.

4. The formula for the maximum number of regions defined by n lines in the plane.

8.5 Good Manners with Induction Proofs 149

Exercise 8.5 Identify the flaw in the proof of the following false statements.

1. WRONG THEOREM. For any natural number n, the following holds:

2 + 4 + \textperiodcentered{} \textperiodcentered{} \textperiodcentered{} + 2n = (n $-$1)(n + 2).

WRONG PROOF. Assume that 2 + 4 + \textperiodcentered{} \textperiodcentered{} \textperiodcentered{} + 2k = (k $-$1)(k + 2) for some k $\in$N. Then

2 + 4 + \textperiodcentered{} \textperiodcentered{} \textperiodcentered{} + 2(k + 1) = (2 + 4 + \textperiodcentered{} \textperiodcentered{} \textperiodcentered{} + 2k) + 2(k + 1)

= (k $-$1)(k + 2) + 2(k + 1) (by the induction hypothesis)

= k(k + 3) = ((k + 1) $-$1))((k + 1) + 2)

which completes the induction. 

2. WRONG THEOREM. For any n $\in$N, if max(i, j) = n for two natural numbers i and j, then i = j. Thus any two natural numbers are equal.

WRONG PROOF. The statement is true for n = 1: if max(i, j) = 1, then, necessarily, i = j = 1. Now assume the statement to be true for some k $\geq$1, and let i, j be natural numbers such that max(i, j) = k + 1. Then max(i $-$1, j $-$1) = k, and by the induction hypothesis, we have that i $-$1 = j $-$1. But then i = j, and therefore our statement is true for n = k + 1. This completes the induction. 

3. WRONG THEOREM. (Pólya8) In any group of n horses, all horses have the same colour.

WRONG PROOF. The statement is clearly true for n = 1. Assume now that the statement is true for some n = k $\geq$1, and consider an arbitrary collection of k + 1 horses.

\textbullet{} \textbullet{} \textperiodcentered{} \textperiodcentered{} \textperiodcentered{} \textbullet{} \textbullet{}

k+1

By the induction hypothesis, the first k horses have the same colour, and so do the last k horses.

k

\textbullet{}

\textbullet{} \textbullet{} \textperiodcentered{} \textperiodcentered{} \textperiodcentered{} \textbullet{}

k

But then all k + 1 horses have the same colour as the k $-$1 horses common to the two sets.

\textbullet{} \textbullet{} \textperiodcentered{} \textperiodcentered{} \textperiodcentered{} \textbullet{}

\textbullet{}

k$-$1

This completes the inductive step. 

8 George Pólya (Hungarian: 1887--1985).

\end{document}